\newpage
\hypertarget{para:meravigliosamente}{}

\mypoemtitle{Meravigliosamente - Giacomo da Lentini}{Canzonetta}

\textbf{Commento introduttivo:}\\
In questa canzonetta, Giacomo da Lentini affronta uno dei temi cardinali della poesia Siciliana ovvero l'espressione dell'amore come esperienza visiva, viene applicato inoltre un altro topos letterario tipico della letteratura provenzale ovvero quello dell'amante timido ed incapace di rivelare i propri sentimenti alla donna amata.

\vspace{2em}

\begin{adjustwidth}{2em}{0pt}
\poemlines{5}
\verselinenumbersleft
\begin{verse}
Meravigliosa - mente \label{m1} \\
un amor mi distringe \label{m2} \\
e mi tene ad ogn’ora. \label{m3} \\
Com’om che pone mente \label{m4} \\
in altro exemplo \textbf{pinge} \label{m5} \\
la simile \textbf{pintura}, \label{m6} \\
così, bella, facc’eo, \label{m7} \\
che ’nfra lo core meo \label{m8} \\
porto la tua \textbf{figura}. \label{m9} \\!

In cor par ch’eo vi porti, \label{m10} \\
\textbf{pinta} come parete, \label{m11} \\
e non pare difore. \label{m12} \\
O Deo, co’ mi par forte \label{m13} \\
non so se lo sapete, \label{m14} \\
con’ v’amo di bon core; \label{m15} \\
ch’eo son sì vergognoso \label{m16} \\
ca pur vi guardo ascoso \label{m17} \\
e non vi mostro amore. \label{m18} \\!

Avendo gran disio \label{m19} \\
\textbf{dipinsi una pintura}, \label{m20} \\
bella, voi simigliante, \label{m21} \\
e quando voi non vio \label{m22} \\
guardo ’n quella \textbf{figura}, \label{m23} \\
par ch’eo v’aggia davante: \label{m24} \\
come quello che crede \label{m25} \\
salvarsi per sua fede, \label{m26} \\
ancor non veggia inante. \label{m27} \\!

Al cor m’ard’una doglia, \label{m28} \\
com’ om che ten lo foco \label{m29} \\
a lo suo seno ascoso, \label{m30} \\
e quando più lo ’nvoglia, \label{m31} \\
allora arde più loco \label{m32} \\
e non pò star incluso: \label{m33} \\
similemente eo ardo \label{m34} \\
quando pass’e non guardo \label{m35} \\
a voi, vis’amoroso. \label{m36} \\!

S’eo guardo, quando passo, \label{m37} \\
inver’ voi no mi giro, \label{m38} \\
bella, per risguardare; \label{m39} \\
andando, ad ogni passo \label{m40} \\
getto uno gran sospiro \label{m41} \\
ca facemi ancosciare; \label{m42} \\
e certo bene ancoscio, \label{m43} \\
c’a pena mi conoscio, \label{m44} \\
tanto bella mi pare. \label{m45} \\!

Assai v’aggio laudato, \label{m46} \\
madonna, in tutte parti, \label{m47} \\
di bellezze c’avete. \label{m48} \\
Non so se v’è contato \label{m49} \\
ch’eo lo faccia per arti, \label{m50} \\
che voi pur v’ascondete: \label{m51} \\
sacciatelo per singa \label{m52} \\
zo ch’eo no dico a linga, \label{m53} \\
quando voi mi vedite \label{m54} \\!

Canzonetta novella, \label{m55} \\
va’ canta nova cosa; \label{m56} \\
lèvati da maitino \label{m57} \\
davanti a la più bella, \label{m58} \\
fiore d’ogn’amorosa, \label{m59} \\
bionda più c’auro fino: \label{m60} \\
«Lo vostro amor, ch’è caro, \label{m61} \\
donatelo al Notaro \label{m62} \\
ch’è nato da Lentino». \label{m63}
\end{verse}
\end{adjustwidth}

\vspace{3em}

\begin{commentaries}
    \scolio[\ref{m1}]{Meravigliosamente} L'avverbio definisce lo stupore dell'amante davanti alla bellezza.
    \scolio[m4]{Com'omo che pone mente} Similitudine col pittore che osserva un modello per ritrarlo.
    \scolio[m8]{'nfra lo core meo} L'immagine della donna custodita interiormente.
    \scolio[m11]{pinta come parete} Come se l'immagine fosse affrescata nel cuore.
    \scolio[m25]{come quello che crede} Similitudine religiosa: l'amore come fede in ciò che non si vede direttamente.
    \scolio[m62]{al Notaro} Firma del poeta, identificato come Giacomo da Lentini.
\end{commentaries}